% from ega0-1-fr.tex
\oldpage[0\textsubscript{I}]{74}

de $A\{S^{-1}\}$~; en outre, le corps des fractions de
$A\{S^{-1}\}/\mathfrak p\{S^{-1}\}$ est canoniquement isomorphe à celui de
$A/\mathfrak p$.
\end{proposition}

\begin{proposition}[7.6.11]
\label{0.7.6.11-fr}
\begin{enumerate}
  \item[\rm(i)] Si $A$ est un anneau admissible, il en est de même de
    $A'=A\{S^{-1}\}$ et pour tout idéal de définition $\mathfrak J$ de $A$,
    $\mathfrak J'=\mathfrak J\{S^{-1}\}$ est un idéal de définition de $A'$.
  \item[\rm(ii)] Soient $A$ un anneau adique, $\mathfrak J$ un idéal de
    définition de $A$ tel que $\mathfrak J/\mathfrak J^2$ soit de type fini
    sur $A/\mathfrak J$~; alors $A'$ est un anneau $\mathfrak J'$-adique et
    $\mathfrak J'/\mathfrak J'^2$ est de type fini sur
    $A'/\mathfrak J'$. Si en outre $A$ est noethérien, il en est de même de
    $A'$.
\end{enumerate}
\end{proposition}

\begin{enumerate}
  \item[\rm(i)] Si $\mathfrak J$ est un idéal de définition dans $A$, il est
    clair que $S^{-1}\mathfrak J$ est un idéal de définition dans l'anneau
    topologique $S^{-1}A$, car on a
    $(S^{-1}\mathfrak J)^n=S^{-1}\mathfrak J^n$. Soit $A''$ l'anneau séparé
    associé à $S^{-1}A$, $\mathfrak J''$ l'image de $S^{-1}\mathfrak J$ dans
    $A''$~; l'image de $S^{-1}\mathfrak J^n$ est $\mathfrak J''^n$, donc
    $\mathfrak J''^n$ tend vers 0 dans $A''$~; comme $\mathfrak J'$ est
    l'adhérence de $\mathfrak J''$ dans $A'$, $\mathfrak J'^n$ est contenu
    dans l'adhérence de $\mathfrak J''^n$, donc tend vers 0 dans $A'$.
  \item[\rm(ii)] Posons $A_i=A/\mathfrak J^{i+1}$, et pour $i\leq j$, soit
    $u_{ij}$ l'homomorphisme canonique
    $A/\mathfrak J^{j+1}\to A/\mathfrak J^{i+1}$~; soit $S_i$ l'image
    canonique de $S$ dans $A_i$, et posons $A_i'=S_i^{-1}A_i$~; soit enfin
    $u_{ij}':A_j'\to A_i'$ l'homomorphisme déduit canoniquement de $u_{ij}$.
    Montrons que le système projectif $(A_i',u_{ij}')$ vérifie les conditions
    de la prop.~(7.2.7)~: il est clair que les $u_{ij}'$ sont surjectifs~;
    d'autre part, le noyau de $u_{ij}'$ est
    $S_j^{-1}(\mathfrak J^{i+1}/\mathfrak J^{j+1})$ (1.3.2), égal à
    $\mathfrak J_j'^{i+1}$, où
    $\mathfrak J_j'=S_j^{-1}(\mathfrak J/\mathfrak J^{j+1})$~; enfin,
    $\mathfrak J_1'/\mathfrak J_1'^2=S_1^{-1}(\mathfrak J/\mathfrak J^2)$,
    et comme $\mathfrak J/\mathfrak J^2$ est de type fini sur
    $A/\mathfrak J^2$, $\mathfrak J_1'/\mathfrak J_1'^2$ est de type fini
    sur $A_1'$. Enfin, si $A$ est noethérien, il en est de même de
    $A_0'=S_0^{-1}(A/\mathfrak J)$, ce qui achève de démontrer la proposition
    (7.2.8).
\end{enumerate}

\begin{corollary}[7.6.12]
\label{0.7.6.12-fr}
Sous les hypothèses de (7.6.11, (ii)), on a
$(\mathfrak J\{S^{-1}\})^n=\mathfrak J^n\{S^{-1}\}$.
\end{corollary}

Cela résulte en effet de (7.2.7) et de la démonstration de (7.6.11).

\begin{proposition}[7.6.13]
\label{0.7.6.13-fr}
Soient $A$ un anneau adique noethérien, $S$ une partie multiplicative de $A$~;
alors $A\{S^{-1}\}$ est un $A$-module plat.
\end{proposition}

En effet, si $\mathfrak J$ est un idéal de définition de $A$,
$A\{S^{-1}\}$ est le séparé complété de l'anneau noethérien $S^{-1}A$ muni
de la topologie $S^{-1}\mathfrak J$-préadique~; par suite (7.3.3)
$A\{S^{-1}\}$ est un $S^{-1}A$-module plat~; comme $S^{-1}A$ est un
$A$-module plat (6.3.1), la proposition résulte de la transitivité de la
platitude (6.2.1).

\begin{corollary}[7.6.14]
\label{0.7.6.14-fr}
Sous les hypothèses de (7.6.13), soit $S'\subset S$ une seconde partie
multiplicative de $A$~; alors $A\{S^{-1}\}$ est un
$A\{{S'}^{-1}\}$-module plat.
\end{corollary}

En effet (7.6.8), $A\{S^{-1}\}$ s'identifie canoniquement à
$A\{{S'}^{-1}\}\{S_0^{-1}\}$, où $S_0$ est l'image canonique de $S$ dans
$A\{{S'}^{-1}\}$, et $A\{{S'}^{-1}\}$ est noethérien (7.6.11).

\begin{env}[7.6.15]
\label{0.7.6.15-fr}
Pour tout élément $f$ d'un anneau linéairement topologisé $A$, nous
désignerons par $A_{\{f\}}$ l'anneau complet de fractions $A\{S_f^{-1}\}$,
où $S_f$ est l'ensemble multiplicatif des $f^n$ ($n\geq 0$)~; pour tout
idéal ouvert $\mathfrak a$ de $A$, nous écrirons $\mathfrak a_{\{f\}}$ au
lieu de $\mathfrak a\{S_f^{-1}\}$. Si $g$ est un second élément de $A$, on a
un homomorphisme canonique continu $A_{\{f\}}\to A_{\{fg\}}$ (7.6.7).
Lorsque $f$ parcourt une partie multiplicative $S$ de $A$, les $A_{\{f\}}$
forment donc un système inductif filtrant d'anneaux pour les homomorphismes
précédents~; nous poserons
$A_{\{S\}}=\varinjlim_{f\in S}A_{\{f\}}$. Pour tout $f\in S$, on a un
homomorphisme $A_{\{f\}}\to A\{S^{-1}\}$ (7.6.7), et

